% Generated draft: requires mathematical and source review.
\documentclass[11pt]{article}
\usepackage[margin=25mm]{geometry}
\usepackage{fontspec}
\setmainfont{Latin Modern Roman}
\usepackage{amsmath,amssymb,mathtools,unicode-math}
\setmathfont{Latin Modern Math}
\usepackage{enumitem,needspace,xurl,hyperref}
\setlength{\parindent}{0pt}
\setlength{\parskip}{5pt}
\setlength{\emergencystretch}{3em}
\newcommand{\N}{\mathbb{N}}
\newcommand{\Z}{\mathbb{Z}}
\newcommand{\Q}{\mathbb{Q}}
\newcommand{\R}{\mathbb{R}}
\newcommand{\C}{\mathbb{C}}
\newcommand{\F}{\mathbb{F}}
\newcommand{\A}{\mathbb{A}}
\newcommand{\PP}{\mathbb{P}}
\newcommand{\E}{\mathbb{E}}
\newcommand{\eps}{\varepsilon}
\newcommand{\dd}{\,\mathrm{d}}
\begin{document}
\section*{UnsolvedMath ID: 30000329}
\textbf{Source problem code:} catalogue entry
\subsection{Definitions and mathematical statement}
A limit group, in the standard finitely generated convention, is a finitely generated fully residually free group: for every finite subset \(S\) there is a homomorphism to a free group that is injective on \(S\). A hyperbolic limit group is one whose Cayley graph is Gromov hyperbolic. The first question is whether every such group \(L\) has a compact negatively curved classifying space \(X=K(L,1)\), meaning \(\pi_1(X)\cong L\) and contractible universal cover, with the required negative-curvature geometry.

In a common metric interpretation the requested space is compact and locally \(\operatorname{CAT}(-1)\). The source does not specify this versus another negative-curvature convention. Requiring a closed smooth manifold of negative sectional curvature would be a substantially different assertion and is not imposed. The exact admissible complex and metric class must be checked in the report.

Countable groups \(G,H\) are measure equivalent if there is a nonzero standard \(\sigma\)-finite measure space with commuting, measure-preserving, essentially free actions of \(G\) and \(H\), each admitting a measurable fundamental domain of finite positive measure. The second question asks for a characterization of countable groups measure equivalent to some finitely generated free group \(F_r\). Its particular subquestion is
\[
\forall L\text{ a limit group},\quad\exists r<\infty\quad L\text{ is measure equivalent to }F_r.
\]
The source does not restrict this second question to hyperbolic limit groups. Nor does it specify whether \(r=0\) or \(r=1\) is permitted or only nonabelian free groups \(r\geq2\) are meant. This matters for trivial and abelian limit groups. The displayed finite-rank target retains that ambiguity rather than silently adding nonabelian or hyperbolic hypotheses.

The supplied literature assessment checked 2026-08-21 reports strong geometric and measure-equivalence results for elementarily free groups, meaning groups elementarily equivalent to nonabelian free groups in the intended finitely generated setting, while describing corresponding all-limit-group questions as open or incomplete. No exact theorem statements or hypotheses are supplied, so the two questions are not asserted equivalent or resolved by that contextual statement.
\subsection{Short English statement}
Ask whether hyperbolic limit groups have compact classifying spaces with the intended negative curvature, and separately classify groups measure equivalent to finite-rank free groups, especially all limit groups. The second question includes limit groups beyond the hyperbolic ones.
\subsection{Sources}
\begin{itemize}
\item Local UnsolvedMath registry, permanent ID 30000329, catalogue entry, \emph{Geometry and Measure Equivalence of Limit Groups}.
\item \emph{Low-Dimensional Manifolds} (2006), Oberwolfach Reports, pp. 2519--2570, DOI \url{https://doi.org/10.4171/owr/2005/45}.
\item Supplied literature triage checked 2026-08-21 cites only the original report; no separate geometric or measure-equivalence theorem reference is supplied.
\end{itemize}
\end{document}
